% Standalone solution dossier for prop:a5-block-stability.
\documentclass[../main.tex]{subfiles}
\begin{document}

% === SOLUTION HEADER (proof provenance) =====================================
% ledger-node: prop:a5-block-stability
% refines: prop:a5-block-stability
% bounded_by: obs:symmetry-vs-physical
% evidence_run: none
% checked_by: agent
% author: /root/a4_a5
% reviewer: /root/review_a4_a5
% dossier_editor: /root/review_a4_a5
% review: research/reviews/2026-08-21-a4-a5-independent-agent-audit.md
% date: 2026-08-21
% ============================================================================

\section*{Solution: density-ratio stability of symmetry blocks}

\begin{proposition}
Let $\pi$ and $\widetilde\pi$ be invariant under the same finite isometry group and have the
same pointwise carr\'e du champ. If
\[
 \operatorname*{osc}\log\frac{\dd\widetilde\pi}{\dd\pi}\le\varepsilon,
\]
then, for $b\in\{\mathrm{inv},\mathrm{noninv}\}$,
\[
 e^{-\varepsilon}\lambda_b(\pi)
 \le\lambda_b(\widetilde\pi)
 \le e^{\varepsilon}\lambda_b(\pi).
\]
Consequently an ordering with logarithmic margin greater than $2\varepsilon$ is preserved.
\end{proposition}

\begin{proof}
Set $h=\dd\widetilde\pi/\dd\pi$, $m=\operatorname*{ess\,inf}h$, and
$M=\operatorname*{ess\,sup}h$. Finite log-oscillation gives $M/m\le e^\varepsilon$ and also
makes the two Sobolev domains identical. For every $f$,
\[
 m\Var_\pi(f)\le\Var_{\widetilde\pi}(f)\le M\Var_\pi(f),
 \qquad
 m\mathcal E_\pi(f)\le\mathcal E_{\widetilde\pi}(f)\le M\mathcal E_\pi(f),
\]
where the variance comparison follows from
$\Var_\mu(f)=\inf_c\int(f-c)^2\dd\mu$. Because both measures are invariant, the pointwise
average $P_G$ has the same invariant range and the same kernel for both. Comparing Rayleigh
quotients on either fixed block yields
$(m/M)\lambda_b(\pi)\le\lambda_b(\widetilde\pi)\le(M/m)\lambda_b(\pi)$, proving the claim.
Finally,
$\lambda_{\mathrm{noninv}}(\pi)<e^{-2\varepsilon}\lambda_{\mathrm{inv}}(\pi)$ implies
$\lambda_{\mathrm{noninv}}(\widetilde\pi)<\lambda_{\mathrm{inv}}(\widetilde\pi)$ by applying
opposite sides of the two comparisons.
\end{proof}

\paragraph{Scope.}
This compares block gaps, not individual eigenvectors or spectral projectors.

\end{document}
